% 第三部分 B：讨论（由 main_zh.tex 引入）。

\section{讨论}\label{sec:discussion}

\paragraph{排序好还不够。} 冻结的 \kev{} 对未标注解答的排序很好（AUROC \NUM{fillauc_all_subkev}），用它打分却仍然比丢掉这些解答更差。它在正确性之外所付费的两样东西，后果不同。它的慷慨——$\pc$ 远高于正确解答所占比例——让每个解答的期望奖励提高大致相同的量，会被组基线消掉。它对答完之后还继续写的解答的偏好则不会：这样的文字不改变验证器的标签，所以已标注的解答永远不会惩罚它，而评判模型在未标注的解答上为它付费（\S\ref{sec:res-repeat}）。\texttt{sub\_kev} 的两个策略都学会了重复自己。评价一个补全奖励，要看它在判定相同的解答之间为什么付费；对判定的 AUROC 衡量不了这一点。

\paragraph{重新校准修正了数值，没有修正偏好。} 在迄今的标签上重新拟合输出头，让补全奖励变得校准良好（ECE \NUM{fillece_all_subrecal} 对 \NUM{fillece_all_subkev}），可它的两个策略中仍有一个学会了重复自己。一个看似合理、但这些 run 没有直接检验的原因是：重新拟合的输出头从发布版出发并被锚定在它附近，因此在标签不置可否的地方都保留了发布版输出头的偏好，而答完之后的重复不改变标签，正是这样的地方。从头训练的输出头没有偏好可保留，它的两个策略都没有重复自己。

\paragraph{探针的增益伴随着风格的改变。} \texttt{sub\_probe} 最终在 MATH-500 上比 \texttt{rlvr\_sub} 高 \NUM{d_subprobe_math} 个百分点，在 AIME24 上高 \NUM{d_subprobe_aime24}，也高于满标注的 \texttt{rlvr}。它的两个策略都改成了先写草稿，本文其他 run 都没有；在本项目的其他 GRPO run 中，同样的切换无论奖励为何都标志着最好的那些 run，其中 \NUM{allveronlydraft} 个只用验证器奖励训练（第二部分 \S8.3）。所以每个条件只有两个种子，我们说不清增益中有多少应归功于探针。如果探针确实让切换更容易发生——两个 run 中两个都切换，暗示了这一点但不能证明——那么一个只在四分之一标签上从头训练的输出头，会是比它所基于的评判模型更好的补全奖励，而原因与它预测标签的准确度关系不大。更多 \texttt{sub\_probe} 和 \texttt{rlvr\_sub} 的种子，以及强制或禁止草稿开头的 run，能给出答案。

\paragraph{局限。} 一个模型（4B）、一个任务、一个标注比例（四分之一）、200 步、每个条件两个种子；重新校准的超参数（锚定、缓冲区、步数）只设了一次，没有调。区间按题目重采样，不按种子。风格量是在读过最终解答之后定义的。重复探针用发布版输出头打分，而不是用训练结束时被重新拟合的输出头。训练期间其他实验共用 GPU，所以我们不报告训练时间。
